\section{Remarks}

In this paper we have considered learning as the process of deducing a program for performing a task, from information that does not provide an explicit description of such a program. We have given precise meaning to this notion of learning and have shown that in some restricted but nontrivial contexts it is computationally feasible.

Consider a world containing robots and elephants. Suppose that one of the robots has discovered a recognition algorithm for elephants that can be meaningfully expressed in $k$-conjunctive normal form. Our Theorem A implies that this robot can communicate its algorithm to the rest of the robot population by simply exclaiming ``elephant'' whenever one appears.

An important aspect of our approach, if cast in its greatest generality, is that we require the recognition algorithms of the teacher and learner to agree on an overwhelming fraction of only the natural inputs. Their behavior on unnatural inputs is irrelevant and hence descriptions of all possible worlds are not necessary. If followed to its conclusion this idea has considerable philosophical implications: A learnable concept is nothing more than a short program that distinguishes some natural inputs from some others. If such a concept is passed on among a population in a distributed manner substantial variations in meaning may arise. More importantly, what consensus there is will only be meaningful for natural inputs. The behavior of an individual's program for unnatural inputs has no relevance. Hence thought experiments and logical arguments involving unnatural hypothetical situations may be meaningless activities.

The second important aspects of the formulation is that the notion of oracles makes it possible to discuss a whole range of teacher-learner interactions beyond the mere identification of examples. This is significant in the context of artificial intelligence where a human may be willing to go to great lengths to convey his skills to a machine while being unable to articulate the algorithms he himself uses in the practice of the skills. We expect that some explicit programming does become essential for transmitting skills that are beyond certain limits of difficulty. The identification of these limits is a major goal of the line of work proposed in this paper.
